\documentclass[10pt,a4paper]{amsart}
\usepackage[margin=19mm]{geometry}
\usepackage{amsmath,amssymb,mathtools}
\usepackage[scr=rsfs,cal=euler]{mathalfa}
\usepackage{xfrac}
\usepackage[unicode,hidelinks,bookmarks=false]{hyperref}
\newcommand{\ie}{i.e.\ }
\newcommand{\cf}{cf.\ }
\newcommand{\resp}{respectively\ }
\newcommand{\Subsection}{\subsection}
\providecommand{\nobreakdash}{\nobreak\hbox{-}}
\setlength{\emergencystretch}{2em}
\multlinegap0pt
\usepackage{bm}
\usepackage{bbm}
\usepackage{dsfont}
\usepackage{xspace}
\usepackage{cancel}
\usepackage{mathtools}
\usepackage{amsthm}
\makeatletter
\@ifundefined{theorem}{\newtheorem{theorem}{Theorem}}{}
\@ifundefined{lemma}{\newtheorem{lemma}[theorem]{Lemma}}{}
\makeatother
\title[Bounded ratios for Lorentzian matrices]{Bounded ratios for Lorentzian matrices}
\author{Daoji Huang, June Huh, Daniel Soskin,, Botong Wang}
\date{Source: arXiv:2510.25030v2 (CC BY 4.0)}
\begin{document}
\maketitle

\noindent\emph{Source and licence.} Bounded ratios for Lorentzian matrices. Version arXiv:2510.25030v2 (CC BY 4.0), distributed under the Creative Commons Attribution 4.0 International licence, as recorded in the arXiv licence metadata for this exact version and shown on the abstract page https://arxiv.org/abs/2510.25030v2. Full rights notice: licenses/arxiv-2510.25030-CC-BY-4.0.txt.\par\smallskip

\noindent\textit{Editorial scope.} Counted unit: the Lemma whose source label is \texttt{lem:factors} in the article (source0.tex lines 602--637, source label \texttt{lem:factors}), delivered together with the article's own complete original proof of it (source0.tex lines 638--679, one \texttt{proof} environment of 42 lines). No mathematics is written, repaired or re-derived: statement and proof are the article's own text line for line. Required context retained uncounted: none. Every definition, display and label of those lines is retained unchanged. Omitted from this package and not redistributed: 1--601, 680--1673. Nothing inside a retained excerpt is omitted or altered: every retained body is delivered exactly as the article's lines are written, so no omission receipt of any kind accompanies this package. Citations: the retained excerpts carry no citation command at all, so no bibliography entry of the article is transcribed and no reference is redistributed. Reference adaptation: none was needed and none was made; every reference inside the delivered unit resolves inside it, so no unresolved reference or citation is produced. Printed slips kept literal: no printed word, symbol, hypothesis or number of the counted statement or of its proof was altered. Layout and class: the article's own class is \texttt{amsart} with packages including amsmath, amssymb, amsthm, amsfonts, mathrsfs, mathtools, xy, eucal, mathabx, hyperref, tikz, epic, eepic, yfonts; the wrapper is the lane's pinned \texttt{amsart} editorial wrapper at 10pt on A4, which restates the theorem-like environments the retained text needs and adds the article's own macro definitions that the retained text needs (none), with extra wrapper packages (bm, bbm, dsfont, xspace, cancel, mathtools). The delivered numbering is the unit's own. Assets: no retained span contains a figure environment, a graphics-inclusion command, a PDF-page inclusion, a table or a TikZ picture; no image file is copied into this package, the package declares no asset dependency and no image file is redistributed with it. Licence: version arXiv:2510.25030v2 of this article is distributed under the Creative Commons Attribution 4.0 International licence, as recorded in the arXiv licence metadata for this exact version and shown on the abstract page https://arxiv.org/abs/2510.25030v2. A full rights notice is delivered at licenses/arxiv-2510.25030-CC-BY-4.0.txt. Characters: every retained excerpt is pure ASCII, so the delivered engine, XeLaTeX with its default Latin Modern text font, reports no missing character.
\begin{lemma}\label{lem:factors}
Suppose that $A$ is a symmetric matrix of the form
\[
A=\begin{pmatrix}
0 & 1 & 1 & p_{14} & p_{15} \\
1 & 0 & 1 & p_{24} & p_{25} \\
1 & 1 & 0 & p_{34} & p_{35} \\
p_{14} & p_{24} & p_{34} & p_{44} & p_{45} \\
p_{15} & p_{25} & p_{35} & p_{45} & p_{55} \\
\end{pmatrix}.
\]
Then there exist unique matrices $B$ and $C$ such that $A=BCB^{T}$, where 
\[
BCB^{T}=\begin{pmatrix}
1 & 0 & 0 & 0 & 0 \\
0 & 1 & 0 & 0 & 0 \\
0 & 0 & 1 & 0 & 0 \\
b_{14} & b_{24} & b_{34} & 1 & 0 \\
b_{15} & b_{25} & b_{35} & 0 & 1 \\
\end{pmatrix}
\begin{pmatrix}
0 & 1 & 1 & 0 & 0 \\
1 & 0 & 1 & 0 & 0 \\
1 & 1 & 0 & 0 & 0\\
0 & 0 & 0 & c_{44} & c_{45} \\
0 & 0 & 0 & c_{45} & c_{55} \\
\end{pmatrix}
\begin{pmatrix}
1 & 0 & 0 & b_{14} & b_{15} \\
0 & 1 & 0 & b_{24} & b_{25} \\
0 & 0 & 1 & b_{34} & b_{35} \\
0 & 0 & 0 & 1 & 0 \\
0 & 0 & 0 & 0 & 1 \\
\end{pmatrix}.
\]
\end{lemma}

\begin{proof}
For any $p_{14}, p_{15}, p_{24}, p_{25}, p_{34}, p_{35}$, there exist unique  $b_{14}, b_{15}, b_{24}, b_{25}, b_{34}, b_{35}$ such that
%\begin{equation}\label{eq:small matrix}
\[
\begin{pmatrix}
b_{14} & b_{24} & b_{34}\\
b_{15} & b_{25} & b_{35}\\
\end{pmatrix}
\begin{pmatrix}
0 & 1 & 1 \\
1 & 0 & 1 \\
1 & 1 & 0 \\
\end{pmatrix}
=\begin{pmatrix}
-p_{14} & -p_{24} & -p_{34} \\
-p_{15} & -p_{25} & -p_{35} \\
\end{pmatrix}.
\]
This defines the invertible lower triangular matrix $B$.
The matrix $C$ is uniquely determined by the condition $C=B^{-1}A(B^{T})^{-1}$, and it is straightforward to check that $C$ has the required block diagonal shape.
%\[
%C \coloneq B^{-1}A(B^{T})^{-1}=
%\begin{pmatrix}
%0 & 1 & 1 & 0 & 0 \\
%1 & 0 & 1 & 0 & 0 \\
%1 & 1 & 0 & 0 & 0\\
%0 & 0 & 0 & c_{44} & c_{45} \\
%0 & 0 & 0 & c_{45} & c_{55} \\
%\end{pmatrix}. \qedhere
%\]
%for some $c_{44}, c_{45}, c_{55}$, which proves the existence. For the uniqueness, the given form of $BCB^T$ in the lemma implies that $b_{14}, \dots, b_{35}$ must satisfy equation \eqref{eq:small matrix}. Hence, $B$ is uniquely determined. Then $C=B^{-1}A(B^{T})^{-1}$ is also uniquely determined. 
%We define for $i\neq j$ an elementary $5\times 5$ matrix $E_{ij}(\lambda)$ with  ones on the  diagonal,  $\lambda$ at position $(i,j)$, and zeros everywhere else. Then for a $5\times 5$ matrix $X$, the product $E_{ij}(\lambda) \cdot X$ corresponds to the row operation that adds $\lambda$ times row $j$ to row $i$ of $X$. Similarly, the product $X \cdot E_{ij}(\lambda)^T = X \cdot E_{ji}(\lambda)$ corresponds to the column operation that adds $\lambda$ times column $j$ to column $i$ of $X$.
%We notice that
%\[E_{43}\left(\frac{-p_{14}}{2}\right)\cdot E_{42}\left(\frac{-p_{14}}{2}\right) \cdot E_{41}\left(\frac{p_{14}}{2}\right) \cdot A=\begin{pmatrix}
%0 & 1 & 1 & p_{14} & p_{15} \\
%1 & 0 & 1 & p_{24} & p_{25} \\
%1 & 1 & 0 & p_{34} & p_{35} \\
%0 & p_{24} & p_{34} & * & * \\
%p_{15} & p_{25} & p_{35} & p_{45} & p_{55} \\
%\end{pmatrix}\]
%One can eliminate the symmetric $p_{14}$ entry of $A$ by multiplying $A$ on the right by the  transpose of those three elementary matrices above in reversed order. In a similar way, we can eliminate the remaining entries $p_{24},~p_{34},~p_{15},~p_{25},~p_{35}$. Notice that product of several $E_{ij}$ with $i>j$ is a unit lower triangular matrix, and product of several $E_{ij}$ with $i<j$ is a unit upper triangular matrix. We recall that $\left(XY\right)^{T}=Y^TX^T$ for any $X,Y$ and thus $A=BCB^{T}$ as in the claim of the lemma.
\end{proof}

\end{document}
